\documentclass[10pt,aspectratio=169]{beamer}
\usepackage[utf8]{inputenc}
\usepackage[T1]{fontenc}
\usepackage[french]{babel}
\usepackage{amsmath}
\usepackage{amssymb}
\usepackage{booktabs}
\usepackage{tikz}
\usetikzlibrary{shapes,arrows,positioning,calc,decorations.pathreplacing}
\usepackage{array}
\usepackage{multirow}
\usepackage{colortbl}
\usepackage{listings}
\usepackage{graphicx}
\usepackage{pgfplots}
\pgfplotsset{compat=1.17}

% Configuration listings pour code Python
\lstset{
    basicstyle=\ttfamily\footnotesize,
    breaklines=true,
    frame=single,
    backgroundcolor=\color{backcolour},
    commentstyle=\color{codegreen},
    keywordstyle=\color{blue},
    numberstyle=\tiny\color{codegray},
    stringstyle=\color{codepurple},
    language=Python
}

% Thème
\usetheme{Madrid}
\usecolortheme{default}

% Couleurs personnalisées
\definecolor{bleuactu}{RGB}{0,51,102}
\definecolor{vertactu}{RGB}{0,102,51}
\definecolor{rougeactu}{RGB}{153,0,0}
\definecolor{orangeactu}{RGB}{230,126,34}
\definecolor{codegreen}{rgb}{0,0.6,0}
\definecolor{codegray}{rgb}{0.5,0.5,0.5}
\definecolor{codepurple}{rgb}{0.58,0,0.82}
\definecolor{backcolour}{rgb}{0.95,0.95,0.92}

\setbeamercolor{structure}{fg=bleuactu}
\setbeamercolor{frametitle}{bg=bleuactu,fg=white}

% Commandes personnalisées
\newcommand{\highlight}[1]{\textbf{\color{rougeactu}#1}}
\newcommand{\BE}{\text{BE}}
\newcommand{\PM}{\text{PM}}
\newcommand{\TMG}{\text{TMG}}
\newcommand{\PB}{\text{PB}}
\newcommand{\EUR}{€}

% Informations du document
\title{Best Estimate Vie}
\subtitle{Épargne et Prévoyance sous Solvabilité II — Séance 5}
\author[F. Borel-Mathurin]{Fabrice BOREL-MATHURIN}
\institute[ACPR-DEAR et SAF Lyon 1] {
    ACPR - DEAR\\
    et\\
    Laboratoire SAF Université Claude Bernard Lyon 1
}
\date{Novembre 2025}

\usepackage{multicol}
% --- Compilation pdfLaTeX sans erreur (ajouté à la mise en dépôt, voir README « LaTeX ») -------------
% Caractères absents de l'encodage T1, et accents dans les blocs de code (listings ne lit pas l'UTF-8).
\DeclareUnicodeCharacter{2248}{\ensuremath{\approx}}
\DeclareUnicodeCharacter{2264}{\ensuremath{\leq}}
\DeclareUnicodeCharacter{2265}{\ensuremath{\geq}}
\DeclareUnicodeCharacter{2B50}{\ensuremath{\star}}
\DeclareUnicodeCharacter{274C}{\ensuremath{\times}}
\lstset{literate=
  {é}{{\'e}}1 {è}{{\`e}}1 {ê}{{\^e}}1 {ë}{{\"e}}1 {à}{{\`a}}1 {â}{{\^a}}1 {î}{{\^i}}1 {ï}{{\"i}}1
  {ô}{{\^o}}1 {ù}{{\`u}}1 {û}{{\^u}}1 {ü}{{\"u}}1 {ç}{{\c c}}1 {œ}{{\oe}}1 {É}{{\'E}}1 {È}{{\`E}}1
  {À}{{\`A}}1 {Â}{{\^A}}1 {Ç}{{\c C}}1 {€}{{\texteuro}}1 {°}{{\textdegree}}1 {×}{{$\times$}}1
  {≤}{{$\leq$}}1 {≥}{{$\geq$}}1 {≈}{{$\approx$}}1 {→}{{$\rightarrow$}}1 {—}{{---}}1 {«}{{\guillemotleft}}1
  {»}{{\guillemotright}}1}
% ------------------------------------------------------------------------------------------------------

\begin{document}

\frame{\titlepage}

\begin{frame}
\frametitle{Plan de la séance (4h)}
\tableofcontents
\end{frame}

% ============================================
% SECTION 1 : INTRODUCTION
% ============================================
\section{Introduction : Spécificités du BE Vie}

\begin{frame}{Objectifs pédagogiques}
\begin{block}{À l'issue de cette séance, vous serez capable de :}
\begin{enumerate}
    \item Modéliser rigoureusement les contrats d'assurance-vie (épargne, prévoyance, UC)
    \item Projeter les flux futurs de trésorerie sur plusieurs décennies
    \item Maîtriser les hypothèses actuarielles (mortalité, rachat, frais, participation bénéfices)
    \item Choisir entre approche police-par-police et model points
    \item Implémenter ces projections en Python (transition depuis Excel)
    \item Calculer le Best Estimate conforme Solvency II
\end{enumerate}
\end{block}

\vspace{0.3cm}
\begin{alertblock}{Durée de la séance}
\textbf{Cours magistral :} 2h | \textbf{Travaux dirigés :} 2h (Excel + Python)
\end{alertblock}
\end{frame}

\begin{frame}{Cadre réglementaire Solvency II}
\begin{block}{Base réglementaire}
\textbf{Article 77 Directive 2009/138/CE :}\\
« Les provisions techniques correspondent au montant actuel que les entreprises d'assurance devraient payer si elles \highlight{transféraient immédiatement} leurs droits et obligations à une autre entreprise »
\end{block}

\vspace{0.3cm}
\begin{columns}[T]
\begin{column}{0.48\textwidth}
\textbf{Références complémentaires :}
\begin{itemize}
    \item Article 75 : valorisation market-consistent
    \item EIOPA Guidelines on Valuation of Technical Provisions (2015)
    \item ACPR Position-Recommandation 2013-P-01
\end{itemize}
\end{column}
\begin{column}{0.48\textwidth}
\textbf{Décomposition S2 :}
\[
\boxed{\text{PT} = \BE + \text{RM}}
\]
\begin{itemize}
    \item $\BE$ = Best Estimate (valeur actuelle probable)
    \item $\text{RM}$ = Risk Margin (coût du capital)
\end{itemize}
\end{column}
\end{columns}
\end{frame}

\begin{frame}{BE Vie vs BE Non-Vie : différences fondamentales}
\begin{table}
\centering
\small
\begin{tabular}{l|p{5.5cm}|p{5.5cm}}
\toprule
\textbf{Critère} & \textbf{Non-Vie (Séance 3)} & \textbf{Vie (Séance 4)} \\
\midrule
\rowcolor{bleuactu!10}
\textbf{Objet} & Sinistres survenus & Contrats en cours \\
\textbf{Approche} & Triangle de liquidation & Projection de flux futurs \\
\rowcolor{bleuactu!10}
\textbf{Horizon} & 5-15 ans & 20-50 ans \\
\textbf{Risques} & Développement, réouvertures & Biométrique + financier \\
\rowcolor{bleuactu!10}
\textbf{Méthode} & Chain-Ladder, Bornhuetter-Ferguson & Projection ALM \\
\textbf{Paradigme} & Passé → Futur (run-off) & Présent → Futur (prospectif) \\
\rowcolor{bleuactu!10}
\textbf{Actuaire} & « Archéologue » & « Prospectiviste » \\
\bottomrule
\end{tabular}
\end{table}

\vspace{0.3cm}
\begin{exampleblock}{Changement de logique}
En Non-Vie, on \textbf{observe} le développement passé. En Vie, on \textbf{projette} les comportements futurs.
\end{exampleblock}
\end{frame}

\begin{frame}{Formule générale du Best Estimate Vie}
\begin{block}{Définition formelle}
Le Best Estimate à la date $t=0$ est la valeur actuelle probable des flux futurs :
\[
\boxed{\BE_0 = \sum_{t=1}^{T} \mathbb{E}\left[\frac{CF_t}{(1+r_t)^t}\right]}
\]
\end{block}

\textbf{Où :}
\begin{itemize}
    \item $CF_t$ = Cash-flows nets de l'année $t$ (flux sortants - flux entrants)
    \item $r_t$ = Taux d'actualisation sans risque à maturité $t$ (courbe EIOPA)
    \item $T$ = Horizon jusqu'à extinction complète du portefeuille
    \item $\mathbb{E}[\cdot]$ = Espérance sous probabilité risque-neutre
\end{itemize}

\vspace{0.3cm}
\begin{alertblock}{Particularité du BE Vie}
Contrairement au BE Non-Vie, le BE Vie peut être \highlight{négatif} si les primes futures dépassent les engagements (contrats rentables en run-off).
\end{alertblock}

\textbf{Référence :} Règlement Délégué 2015/35, Articles 28-37
\end{frame}

\begin{frame}{Décomposition détaillée des cash-flows}
\textbf{Cash-flows sortants (année $t$) :}
\begin{align*}
CF_t^{\text{out}} = &\underbrace{\text{Prestations décès}_t}_{\text{capitaux + majorations}} 
+ \underbrace{\text{Prestations vie}_t}_{\text{échéances programmées}} \\
&+ \underbrace{\text{Rachats}_t}_{\text{totaux + partiels}} 
+ \underbrace{\text{Arrérages}_t}_{\text{rentes viagères}} \\
&+ \underbrace{\text{Frais}_t}_{\text{acquisition + gestion + inflation}} 
+ \underbrace{\PB_t}_{\text{participation bénéfices}}
\end{align*}

\vspace{0.2cm}
\textbf{Cash-flows entrants (année $t$) :}
\[
CF_t^{\text{in}} = \underbrace{\text{Primes futures}_t}_{\text{périodiques uniquement}} 
+ \underbrace{\text{Pénalités rachat}_t}_{\text{si applicables}}
\]

\vspace{0.2cm}
\begin{exampleblock}{Flux net à actualiser}
\[
\boxed{CF_t = CF_t^{\text{out}} - CF_t^{\text{in}}}
\]
\end{exampleblock}
\end{frame}

% ============================================
% SECTION 2 : TYPOLOGIE DES PRODUITS VIE
% ============================================
\section{Typologie des Produits Vie}

\begin{frame}{Vue d'ensemble des produits Vie}
\begin{center}
\begin{tikzpicture}[
    main/.style={rectangle, draw, fill=bleuactu!30, text width=3cm, text centered, minimum height=1cm, font=\small\bfseries},
    prod/.style={rectangle, draw, fill=vertactu!20, text width=2.5cm, text centered, minimum height=0.8cm, font=\scriptsize}
]
    \node[main] (vie) at (0,0) {Assurance\\Vie};
    
    \node[prod] (euros) at (-4.5,-2) {Contrats\\en euros};
    \node[prod] (uc) at (-1.5,-2) {Contrats\\UC};
    \node[prod] (prev) at (1.5,-2) {Prévoyance\\décès/invalidité};
    \node[prod] (rente) at (4.5,-2) {Rentes\\viagères};
    
    \draw[->, thick] (vie) -- (euros);
    \draw[->, thick] (vie) -- (uc);
    \draw[->, thick] (vie) -- (prev);
    \draw[->, thick] (vie) -- (rente);
    
    \node[below=0.3cm of euros, text width=2.5cm, align=center, font=\tiny] {TMG, PB\\Garantie capital};
    \node[below=0.3cm of uc, text width=2.5cm, align=center, font=\tiny] {Support OPCVM\\Risque assuré};
    \node[below=0.3cm of prev, text width=2.5cm, align=center, font=\tiny] {Capital\\temporaire};
    \node[below=0.3cm of rente, text width=2.5cm, align=center, font=\tiny] {Viager\\Longévité};
\end{tikzpicture}
\end{center}

\vspace{0.3cm}
\textbf{Chaque produit nécessite une modélisation spécifique du Best Estimate}
\end{frame}

\begin{frame}{I. Contrats en euros (fonds général)}
\begin{block}{Définition}
Contrat d'épargne où l'assureur \highlight{garantit le capital} investi et verse une rémunération annuelle minimale (TMG = Taux Minimum Garanti).
\end{block}

\begin{columns}[T]
\begin{column}{0.48\textwidth}
\textbf{Caractéristiques principales :}
\begin{itemize}
    \item Garantie capital à 100\%
    \item TMG : 0\% à 3\% (historique)
    \item Participation aux bénéfices (85-90\%)
    \item Effet cliquet (acquis définitifs)
    \item Liquidité (rachat possible)
\end{itemize}
\end{column}
\begin{column}{0.48\textwidth}
\textbf{Flux à modéliser :}
\begin{itemize}
    \item[$+$] Primes périodiques
    \item[$-$] Prestations décès
    \item[$-$] Rachats (totaux/partiels)
    \item[$-$] Frais de gestion
    \item[$-$] Participation aux bénéfices
\end{itemize}
\end{column}
\end{columns}

\vspace{0.3cm}
\begin{alertblock}{Risque pour l'assureur}
Si $\text{rendement actifs} < \TMG$, l'assureur doit puiser dans ses fonds propres pour garantir le TMG !
\end{alertblock}
\end{frame}

\begin{frame}{Mécanisme de la participation aux bénéfices}
\begin{block}{Principe réglementaire}
\textbf{Article A331-3 Code des Assurances :} Minimum \highlight{85\%} (pratique : 90\%) des produits financiers redistribués aux assurés sur \highlight{8 ans glissants}.
\end{block}

\begin{center}
\begin{tikzpicture}[scale=0.9]
\node[draw, rectangle, fill=bleuactu!20, text width=3cm, align=center] (actifs) at (0,2) {Actifs\\Fonds Général};
\node[draw, rectangle, fill=vertactu!20, text width=3cm, align=center] (prod) at (0,0.5) {Produits\\Financiers\\(100\%)};
\node[draw, rectangle, fill=orangeactu!30, text width=2.5cm, align=center] (pb) at (-2.5,-1.5) {PB Assurés\\(85-90\%)};
\node[draw, rectangle, fill=rougeactu!20, text width=2.5cm, align=center] (ass) at (2.5,-1.5) {Marge\\Assureur\\(10-15\%)};
\node[draw, rectangle, fill=bleuactu!10, text width=2cm, align=center] (ppe) at (0,-3) {PPE\\(lissage)};

\draw[->, thick] (actifs) -- (prod);
\draw[->, thick] (prod) -- (pb);
\draw[->, thick] (prod) -- (ass);
\draw[<->, dashed] (pb) -- (ppe);
\end{tikzpicture}
\end{center}

\textbf{PPE} = Provision pour Participation aux Excédents (réserve de lissage)
\end{frame}

\begin{frame}{Formules de projection : contrat euros}
\textbf{Notation :}
\begin{multicols}{2}
\begin{itemize}
    \item $\PM_t$ : Provision mathématique début année $t$
    \item $P_t$ : Prime annuelle
    \item $q_x^d$ : Probabilité décès à l'âge $x$
    \item $q_t^r$ : Probabilité rachat année $t$
    \item $i_t$ : Taux revalorisation ($\TMG + \PB$)
    \item $f$ : Taux de frais sur PM
    \item $N_t$ : Nombre de têtes en vie
\end{itemize}
\end{multicols}

\vspace{0.2cm}
\begin{block}{Système de projection}
\begin{align}
\text{Décès}_t &= N_t \times \overline{\PM}_t \times q_{x+t-1}^d \\
\text{Rachats}_t &= N_t \times \overline{\PM}_t \times q_t^r \\
\text{Frais}_t &= N_t \times \overline{\PM}_t \times f \times (1+\pi)^{t-1} \\
N_{t+1} &= N_t \times (1 - q_{x+t-1}^d) \times (1 - q_t^r) \\
\overline{\PM}_{t+1} &= \frac{(\overline{\PM}_t + P_t - \overline{\PM}_t \times f) \times (1+i_t)}{(1 - q_{x+t-1}^d)(1 - q_t^r)}
\end{align}
\end{block}

$\pi$ = inflation des frais (typiquement 2\% par an)
\end{frame}

\begin{frame}{Exemple numérique : projection année 1}
\textbf{Données contrat :}
\begin{itemize}
    \item Assuré : homme 40 ans, $\PM_0 = 50\,000$ \EUR, prime annuelle $P = 2\,000$ \EUR
    \item $\TMG = 1.5\%$, frais $f = 0.8\%$, $q_{40}^d = 0.123\%$, $q_1^r = 10\%$
    \item Inflation frais : $\pi = 2\%$
\end{itemize}

\begin{exampleblock}{Calculs année 1}
\begin{align*}
\text{Décès}_1 &= 50\,000 \times 0.00123 = \highlight{61.50\, \EUR} \\
\text{Rachats}_1 &= 50\,000 \times 0.10 = \highlight{5\,000\, \EUR} \\
\text{Frais}_1 &= 50\,000 \times 0.008 \times 1.02^0 = \highlight{400\, \EUR} \\
\text{Flux sortants}_1 &= 61.50 + 5\,000 + 400 = \highlight{5\,461.50\, \EUR} \\
\text{Flux entrants}_1 &= 2\,000\, \EUR \\
\text{Flux net}_1 &= 5\,461.50 - 2\,000 = \highlight{3\,461.50\, \EUR}
\end{align*}
\end{exampleblock}

\textbf{PM fin année 1 :}
\[
\PM_2 = (50\,000 + 2\,000 - 61.50 - 5\,000 - 400) \times 1.015 = \boxed{47\,236.58\, \EUR}
\]
\end{frame}

\begin{frame}{II. Contrats en Unités de Compte (UC)}
\begin{block}{Définition}
Contrat où l'épargne est investie sur des supports financiers (OPCVM, actions) choisis par l'assuré. L'assureur garantit un \highlight{nombre d'UC}, pas une valeur en euros → \textbf{risque transféré à l'assuré}.
\end{block}

\begin{columns}[T]
\begin{column}{0.48\textwidth}
\textbf{Différences vs contrats euros :}
\begin{itemize}
    \item \textcolor{rougeactu}{Pas de garantie capital}
    \item Pas de TMG ni PB
    \item Risque de marché → assuré
    \item Frais plus élevés (1-2\%)
    \item Transparence des supports
\end{itemize}
\end{column}
\begin{column}{0.48\textwidth}
\textbf{Garanties optionnelles :}
\begin{enumerate}
    \item \textbf{Plancher au décès}\\
    {\scriptsize Ex: min 100\% versements}
    \item \textbf{Cliquet périodique}\\
    {\scriptsize Tous les X ans}
    \item \textbf{Garantie à l'échéance}\\
    {\scriptsize Ex: à 8 ans, 100\% versements}
    \item \textbf{Stop-loss dynamique}\\
    {\scriptsize Sécurisation progressive}
\end{enumerate}
\end{column}
\end{columns}

\vspace{0.3cm}
\begin{alertblock}{Risque réintroduit}
Les garanties financières réintroduisent du \textbf{risque de marché} pour l'assureur, qui doit le provisionner dans le BE !
\end{alertblock}
\end{frame}

\begin{frame}{Modélisation UC avec garantie plancher}
\textbf{Notation :}
\begin{itemize}
    \item $V_t$ = Valeur des UC en année $t$
    \item $G$ = Plancher garanti (ex: 100\% versements nets de frais)
    \item $r_t^{UC}$ = Rendement du support UC en année $t$
    \item $f^{UC}$ = Frais de gestion UC (annuels)
\end{itemize}

\vspace{0.3cm}
\begin{block}{Évolution valeur UC}
\[
\boxed{V_{t+1} = V_t \times (1 - f^{UC}) \times (1 + r_t^{UC})}
\]
\end{block}

\begin{block}{Coût de la garantie au décès (année $t$)}
\[
\boxed{\text{Coût garantie}_t = \max(0, \, G - V_t) \times q_{x+t-1}^d \times N_t}
\]
\end{block}

\begin{block}{Best Estimate de la garantie}
\[
\boxed{\BE^{\text{garantie}} = \sum_{t=1}^{T} \frac{\text{Coût garantie}_t}{(1+r_t)^t}}
\]
\end{block}
\end{frame}

\begin{frame}{Exemple : UC avec garantie plancher en crise}
\textbf{Scénario de crise :} $V_0 = 100\,000$ \EUR, $G = 100\,000$ \EUR, $f^{UC} = 1\%$, $r^{UC} = -3\%$ par an

\vspace{0.3cm}
\begin{center}
\begin{tikzpicture}[scale=0.75]
\begin{axis}[
    xlabel={Année},
    ylabel={Valeur (\EUR)},
    grid=major,
    legend pos=south west,
    width=11cm,
    height=6cm,
    ymin=65000,
    ymax=105000
]
\addplot[color=bleuactu, thick, mark=*] coordinates {
    (0, 100000) (1, 96030) (2, 92285) (3, 88760) (4, 85449)
    (5, 82344) (6, 79439) (7, 76727)
};
\addplot[color=rougeactu, thick, dashed, line width=1.5pt] coordinates {
    (0, 100000) (7, 100000)
};
\legend{Valeur UC (baisse 3\% + frais 1\%), Plancher garanti}
\end{axis}
\end{tikzpicture}
\end{center}

\textbf{Interprétation :} Après 7 ans, valeur UC = 76 727 \EUR < plancher 100 000 \EUR.\\
Si décès en année 7 : garantie activée → \highlight{coût = 23 273 \EUR} pour l'assureur !
\end{frame}

\begin{frame}{III. Contrats de Prévoyance}
\begin{block}{Définition}
Contrats couvrant les \highlight{risques biométriques} : décès, incapacité, invalidité, dépendance.
\end{block}

\begin{table}
\centering
\small
\begin{tabular}{llp{6cm}}
\toprule
\textbf{Produit} & \textbf{Garantie} & \textbf{Caractéristiques} \\
\midrule
\rowcolor{bleuactu!10}
Décès temporaire & Capital & Si décès avant échéance (ex: prêt immobilier) \\
Décès vie entière & Capital & Quel que soit l'âge au décès \\
\rowcolor{bleuactu!10}
Invalidité/IPT & Rente & En cas d'Incapacité Permanente Totale \\
Dépendance & Rente & Si perte d'autonomie (critères ADL) \\
\bottomrule
\end{tabular}
\end{table}

\vspace{0.3cm}
\begin{exampleblock}{Formule BE décès temporaire (durée $n$ ans)}
\[
\boxed{\BE = \sum_{t=1}^{n} \frac{K \times {}_{t-1}p_x \times q_{x+t-1}^d - P \times {}_{t-1}p_x \times (1-l)^{t-1}}{(1+r_t)^t}}
\]

Où : $K$ = capital garanti, $P$ = prime annuelle, $l$ = taux de résiliation, ${}_{t}p_x$ = probabilité survie jusqu'à $t$
\end{exampleblock}
\end{frame}

\begin{frame}{Tables de mortalité et ajustements}
\begin{columns}[T]
\begin{column}{0.48\textwidth}
\textbf{Tables réglementaires françaises :}
\begin{itemize}
    \item \textbf{TH/TF 00-02}\\
    {\scriptsize Tables d'expérience 2000-2002}
    \item \textbf{TGH/TGF 05}\\
    {\scriptsize Tables générationnelles (projections)}
    \item \textbf{TD/TV 88-90}\\
    {\scriptsize Tables pour rentes viagères}
\end{itemize}

\vspace{0.3cm}
\begin{alertblock}{Choix crucial}
\textbf{Épargne/décès :} TH/TF 00-02\\
\textbf{Rentes viagères :} TGH/TGF 05\\
(risque de longévité)
\end{alertblock}
\end{column}
\begin{column}{0.48\textwidth}
\textbf{Ajustement expérience propre :}
\[
\boxed{q_x^{\text{ajustée}} = \alpha \times q_x^{\text{table}}}
\]

\textbf{Valeurs typiques :} $\alpha \in [0.7, \, 1.3]$

\vspace{0.3cm}
\begin{itemize}
    \item $\alpha < 1$ : mortalité plus faible\\
    {\scriptsize (portefeuille aisé, rentes)}
    \item $\alpha > 1$ : mortalité plus forte\\
    {\scriptsize (garanties décès)}
\end{itemize}

\vspace{0.3cm}
\textbf{Justification requise :}
\begin{itemize}
    \item Études statistiques internes
    \item Validation actuaire
    \item Documentation ACPR
\end{itemize}
\end{column}
\end{columns}

\textbf{Référence :} Arrêté du 1er août 2016 (tables réglementaires)
\end{frame}

\begin{frame}{IV. Rentes viagères}
\begin{block}{Définition}
Contrat où l'assureur verse périodiquement une somme jusqu'au \highlight{décès du rentier}.
\end{block}

\begin{columns}[T]
\begin{column}{0.48\textwidth}
\textbf{Types de rentes :}
\begin{itemize}
    \item \textbf{Immédiate :} versements démarrent immédiatement
    \item \textbf{Différée :} phase d'accumulation puis rente
    \item \textbf{Réversible :} continue sur conjoint survivant
    \item \textbf{Annuités garanties :} min X années versées
    \item \textbf{Indexée :} revalorisation annuelle
\end{itemize}
\end{column}
\begin{column}{0.48\textwidth}
\textbf{Formule rente viagère simple :}
\[
\boxed{\BE = R \times \sum_{t=1}^{\omega - x} \frac{{}_{t-1}p_x}{(1+r_t)^t}}
\]

Où :
\begin{itemize}
    \item $R$ = arrérages annuels
    \item $\omega$ = âge limite (120 ans)
    \item ${}_{t}p_x$ = $\prod_{k=0}^{t-1}(1-q_{x+k})$
\end{itemize}
\end{column}
\end{columns}

\vspace{0.3cm}
\begin{exampleblock}{Calcul numérique : homme 65 ans, 12 000 \EUR/an}
\[
\BE = 12\,000 \times \underbrace{\sum_{t=1}^{55} \frac{{}_{t-1}p_{65}}{(1.025)^t}}_{= \ddot{a}_{65} \approx 15.8} \approx 12\,000 \times 15.8 = \boxed{189\,600\, \EUR}
\]
\end{exampleblock}
\end{frame}

\begin{frame}{Risque de longévité : enjeu majeur}
\begin{alertblock}{Risque systématique}
Si les rentiers vivent \textbf{plus longtemps que prévu} dans les tables, les versements dépassent significativement les provisions constituées !
\end{alertblock}

\begin{center}
\begin{tikzpicture}[scale=0.75]
\begin{axis}[
    xlabel={Âge},
    ylabel={Espérance de vie résiduelle (années)},
    grid=major,
    legend pos=north east,
    width=11cm,
    height=6cm
]
\addplot[color=bleuactu, thick, mark=square*] coordinates {
    (65, 19.5) (70, 15.8) (75, 12.5) (80, 9.5) (85, 7.0) (90, 5.0)
};
\addplot[color=rougeactu, thick, dashed, mark=*] coordinates {
    (65, 21.0) (70, 17.0) (75, 13.5) (80, 10.3) (85, 7.6) (90, 5.5)
};
\legend{Prévision initiale (TGH/TGF 05), Réalisation observée (+1.5 ans)}
\end{axis}
\end{tikzpicture}
\end{center}

\textbf{Impact financier :} +1.5 ans d'espérance de vie à 65 ans → \highlight{+7.7\%} sur le BE rentes !

\textbf{Solutions :} réassurance longévité, tables générationnelles, clauses d'indexation limitée

\textbf{Référence :} Cahiers de la Chaire transitions démographiques (2023)
\end{frame}

% ============================================
% SECTION 3 : MODÉLISATION DES FLUX
% ============================================
\section{Modélisation des Flux}

\begin{frame}{Approche police par police vs model points}
\begin{table}
\centering
\small
\begin{tabular}{l|p{5.5cm}|p{5.5cm}}
\toprule
\textbf{Critère} & \textbf{Police par police} & \textbf{Model points} \\
\midrule
\rowcolor{bleuactu!10}
\textbf{Principe} & Projection individuelle de chaque contrat & Regroupement en contrats représentatifs \\
\textbf{Précision} & Maximale (toutes options/garanties) & Approximation contrôlée \\
\rowcolor{bleuactu!10}
\textbf{Temps calcul} & Très élevé (plusieurs heures) & Réduit (×100 à ×1000) \\
\textbf{Données} & Nécessite qualité parfaite & Robuste aux données manquantes \\
\rowcolor{bleuactu!10}
\textbf{Usage} & Portefeuilles hétérogènes, < 10 000 contrats & Production industrielle, > 10 000 contrats \\
\textbf{Validation} & Non nécessaire & Back-testing obligatoire \\
\bottomrule
\end{tabular}
\end{table}

\vspace{0.3cm}
\begin{exampleblock}{Ordre de grandeur}
50 000 contrats réels → 200-500 model points (gain : ×100 à ×250 en rapidité)
\end{exampleblock}
\end{frame}

\begin{frame}{Construction des model points}
\begin{block}{Critères de segmentation typiques}
\begin{itemize}
    \item \textbf{Âge :} [<30], [30-40], [40-50], [50-60], [60-70], [70+]  \hfill (6 classes)
    \item \textbf{Ancienneté :} [0-2], [2-5], [5-10], [10-15], [15+]  \hfill (5 classes)
    \item \textbf{PM :} [<10k], [10k-50k], [50k-100k], [100k-250k], [>250k]  \hfill (5 classes)
    \item \textbf{Type :} Euros pur, Euros+UC, UC pur, Prévoyance  \hfill (4 classes)
    \item \textbf{TMG :} [0\%], [0-1\%], [1-2\%], [2-3\%], [>3\%]  \hfill (5 classes)
\end{itemize}
\end{block}

\textbf{Nombre théorique de MP :} $6 \times 5 \times 5 \times 4 \times 5 = 3\,000$ combinaisons possibles

\vspace{0.3cm}
\begin{alertblock}{En pratique}
\begin{itemize}
    \item Regroupement supplémentaire des classes peu peuplées
    \item Élimination des combinaisons vides (ex: TMG 3\% + âge <30 ans)
    \item \textbf{Résultat final : 200-500 model points} pour un portefeuille de 50 000 contrats
\end{itemize}
\end{alertblock}
\end{frame}

\begin{frame}{Qualité du model pointing : critères de validation}
\begin{block}{1. Homogénéité intra-groupe (à minimiser)}
\[
\boxed{\text{Var}_{\text{intra}} = \frac{1}{n_g} \sum_{i \in \text{groupe } g} (x_i - \bar{x}_g)^2}
\]
Les contrats d'un même MP doivent être similaires.
\end{block}

\begin{block}{2. Hétérogénéité inter-groupes (à maximiser)}
\[
\boxed{\text{Var}_{\text{inter}} = \sum_{g} n_g \times (\bar{x}_g - \bar{x}_{\text{total}})^2}
\]
Les MP doivent être bien différenciés entre eux.
\end{block}

\begin{block}{3. Conservation des agrégats (erreur < 2\%)}
\begin{align*}
\sum_{\text{MP}} \PM_{MP} \times n_{MP} &= \sum_{\text{polices}} \PM_i \quad \text{(test PM)} \\
\left|\sum_{\text{MP}} \BE_{MP} \times n_{MP} - \sum_{\text{polices}} \BE_i\right| &< 2\% \times \sum_{\text{polices}} \BE_i \quad \text{(test BE)}
\end{align*}
\end{block}

\textbf{Validation :} Back-testing sur portefeuille test (20\% des polices)
\end{frame}

\begin{frame}{Hypothèses de projection : lois de rachat}
\begin{block}{Facteurs influençant les rachats}
\begin{multicols}{2}
\begin{itemize}
    \item Ancienneté du contrat
    \item Contexte économique
    \item Fiscalité (seuils 4 et 8 ans)
    \item Montant de la PM
    \item Canal de distribution
    \item Âge de l'assuré
\end{itemize}
\end{multicols}
\end{block}

\begin{columns}[T]
\begin{column}{0.48\textwidth}
\textbf{Loi de rachat statique :}
\[
q_t^r = 
\begin{cases}
15\% & t \in [1, 2] \\
10\% & t \in [3, 4] \\
6\% & t \in [5, 7] \\
3\% & t \geq 8
\end{cases}
\]

\textbf{Effet fiscal :}
\begin{itemize}
    \item Pic avant 4 ans (fiscalité lourde)
    \item Stabilisation après 8 ans (abattement)
\end{itemize}
\end{column}
\begin{column}{0.48\textwidth}
\textbf{Loi de rachat dynamique :}
\[
\boxed{q_t^r = q_t^{r,\text{base}} \times \left(1 + \beta \times \Delta_t\right)}
\]

Où :
\begin{align*}
\Delta_t &= \max\left(0, \frac{r_{\text{marché},t} - r_{\text{contrat},t}}{r_{\text{contrat},t}}\right) \\
\beta &\in [2, 5] \text{ (élasticité)}
\end{align*}

\textbf{Interprétation :}\\
Si taux de marché $\gg$ taux contrat, les rachats explosent !
\end{column}
\end{columns}
\end{frame}

\begin{frame}{Exemple : rachats dynamiques en contexte de remontée des taux}
\textbf{Hypothèses :}
\begin{itemize}
    \item Taux de rachat de base : $q^{r,\text{base}} = 3\%$
    \item Taux contrat (TMG + PB) : $r_{\text{contrat}} = 1.5\%$
    \item Taux de marché (OAT 10 ans) : $r_{\text{marché}} = 4.0\%$
    \item Élasticité : $\beta = 3$
\end{itemize}

\begin{exampleblock}{Calcul du taux de rachat dynamique}
\begin{align*}
\Delta &= \frac{4.0\% - 1.5\%}{1.5\%} = \frac{2.5\%}{1.5\%} = 1.667 \\
q^r &= 3\% \times (1 + 3 \times 1.667) = 3\% \times 6 = \boxed{\highlight{18\%}}
\end{align*}
\end{exampleblock}

\textbf{Impact :} Le taux de rachat est \textbf{multiplié par 6} en contexte de forte remontée des taux !

\vspace{0.3cm}
\begin{alertblock}{Conséquence pour l'assureur}
Rachats massifs → vente forcée d'actifs en moins-value latente → crise de liquidité potentielle (scénario 2022-2023)
\end{alertblock}
\end{frame}

\begin{frame}{Hypothèses de frais}
\begin{block}{Décomposition des frais (3 catégories)}
\textbf{1. Frais d'acquisition $F_A$ (année 1 uniquement) :}
\[
F_A = \alpha_A \times P_1 \quad \text{avec } \alpha_A \in [20\%, 50\%]
\]
Inclut : commercialisation, médical, ouverture dossier

\vspace{0.2cm}
\textbf{2. Frais de gestion $F_{G,t}$ (annuels) :}
\[
F_{G,t} = \alpha_G \times \PM_t \times (1 + \pi)^{t-1} \quad \text{avec } \alpha_G \in [0.5\%, 1.0\%]
\]

\vspace{0.2cm}
\textbf{3. Frais sur versements $F_V$ (chaque prime) :}
\[
F_V = \alpha_V \times P_t \quad \text{avec } \alpha_V \in [2\%, 5\%]
\]
\end{block}

\textbf{Inflation des frais :} $\pi = 2\%$ par an (convention standard Solvency II)

\vspace{0.2cm}
\begin{alertblock}{Impact à long terme}
Sur 30 ans avec inflation 2\% : $(1.02)^{30} = 1.81$ → \highlight{+81\%} sur les frais annuels !
\end{alertblock}
\end{frame}

\begin{frame}{Frais unitaires vs Frais sur encours : distinction pratique}
\begin{block}{Deux natures de frais en assurance-vie}
Dans la réalité industrielle, les frais ne sont pas tous proportionnels à la PM. Il faut distinguer :
\end{block}

\begin{columns}[T]
\begin{column}{0.48\textwidth}
\textbf{Frais sur encours (proportionnels) :}
\[
F_{\text{encours},t} = \alpha_G \times \PM_t \times (1+\pi)^{t-1}
\]

\textbf{Nature :}
\begin{itemize}
    \item Gestion financière
    \item Coûts IT proportionnels
    \item Commissions de gestion
\end{itemize}

\textbf{Évolution :}
\begin{itemize}
    \item Proportionnels à la PM
    \item Inflation appliquée sur $\alpha_G$
\end{itemize}
\end{column}
\begin{column}{0.48\textwidth}
\textbf{Frais unitaires (fixes par contrat) :}
\[
F_{\text{unitaire},t} = F_0 \times (1+\pi)^{t-1} \times N_t
\]

\textbf{Nature :}
\begin{itemize}
    \item Envoi relevés annuels
    \item Traitement administratif
    \item Coûts IT fixes par police
\end{itemize}

\textbf{Évolution :}
\begin{itemize}
    \item Indépendants de la PM
    \item Exemple : 30€/contrat/an
    \item Inflation appliquée directement
\end{itemize}
\end{column}
\end{columns}

\vspace{0.3cm}
\begin{exampleblock}{Formule complète des frais de gestion}
\[
\boxed{F_{G,t}^{\text{total}} = \underbrace{\alpha_G \times \PM_t \times (1+\pi)^{t-1} \times N_t}_{\text{Frais proportionnels}} + \underbrace{F_0 \times (1+\pi)^{t-1} \times N_t}_{\text{Frais unitaires}}}
\]
\end{exampleblock}
\end{frame}

\begin{frame}{Impact des frais unitaires : exemple numérique}
\textbf{Portefeuille de 1000 contrats, PM moyenne = 50 000 €}

\begin{columns}[T]
\begin{column}{0.48\textwidth}
\textbf{Scénario 1 : Gros contrats}
\begin{itemize}
    \item $\alpha_G = 0.8\%$
    \item $F_0 = 30$ €/contrat/an
    \item $\PM_{\text{moy}} = 50\,000$ €
\end{itemize}

\textbf{Calcul année 1 :}
\begin{align*}
F_{\text{prop}} &= 0.008 \times 50\,000 \\
&= 400\, \text{€/contrat} \\
F_{\text{unit}} &= 30\, \text{€/contrat} \\
\text{Poids unit.} &= \frac{30}{430} = \highlight{7\%}
\end{align*}
\end{column}
\begin{column}{0.48\textwidth}
\textbf{Scénario 2 : Petits contrats}
\begin{itemize}
    \item $\alpha_G = 0.8\%$
    \item $F_0 = 30$ €/contrat/an
    \item $\PM_{\text{moy}} = 5\,000$ €
\end{itemize}

\textbf{Calcul année 1 :}
\begin{align*}
F_{\text{prop}} &= 0.008 \times 5\,000 \\
&= 40\, \text{€/contrat} \\
F_{\text{unit}} &= 30\, \text{€/contrat} \\
\text{Poids unit.} &= \frac{30}{70} = \highlight{43\%}
\end{align*}
\end{column}
\end{columns}

\vspace{0.3cm}
\begin{alertblock}{Conséquence pour le BE}
Pour les petits contrats, \textbf{ignorer les frais unitaires} peut sous-estimer significativement le BE ! C'est particulièrement critique pour les portefeuilles de prévoyance (petites primes).
\end{alertblock}

\textbf{Pratique industrielle :} Distinction systématique dans les modèles ALM de production
\end{frame}

\begin{frame}{Participation aux bénéfices : mécanisme détaillé}
\begin{block}{Cadre réglementaire}
\textbf{Article A331-3 Code des Assurances :}
\begin{itemize}
    \item Minimum \highlight{85\%} (pratique : 90\%) des produits financiers redistribués
    \item Sur \highlight{8 ans glissants} (lissage autorisé via PPE)
    \item Clause de PB minimale généralement contractuelle (TMG)
\end{itemize}
\end{block}

\begin{columns}[T]
\begin{column}{0.48\textwidth}
\textbf{Formule simplifiée :}
\[
\boxed{\PB_t = \max\left(\TMG, \, \rho \times (r_{\text{actifs},t} - f_t) \right)}
\]

Où :
\begin{itemize}
    \item $\rho = 0.90$ (taux de redistribution)
    \item $r_{\text{actifs},t}$ = rendement fonds général
    \item $f_t$ = frais financiers
\end{itemize}
\end{column}
\begin{column}{0.48\textwidth}
\textbf{Calcul numérique :}

Hypothèses :
\begin{itemize}
    \item $r_{\text{actifs}} = 3.5\%$
    \item $f = 0.5\%$
    \item $\TMG = 1.0\%$
\end{itemize}

\begin{align*}
\PB &= \max(1.0\%, \, 0.90 \times (3.5\% - 0.5\%)) \\
&= \max(1.0\%, \, 0.90 \times 3.0\%) \\
&= \max(1.0\%, \, 2.7\%) \\
&= \boxed{\highlight{2.7\%}}
\end{align*}
\end{column}
\end{columns}

\vspace{0.3cm}
\textbf{PPE} = Provision pour Participation aux Excédents (max 8 ans de produits financiers)
\end{frame}

\begin{frame}{Courbe d'actualisation EIOPA}
\begin{block}{Principe Solvency II}
Actualisation avec taux \highlight{sans risque} (risk-free rate curve) → Pas de prime de risque de crédit (obligation souveraine AAA uniquement).
\end{block}

\begin{center}
\begin{tikzpicture}[scale=0.75]
\begin{axis}[
    xlabel={Maturité (années)},
    ylabel={Taux (\%)},
    grid=major,
    legend pos=south east,
    width=11cm,
    height=6cm,
    ymin=0, ymax=4
]
\addplot[color=bleuactu, thick, mark=*, line width=1.5pt] coordinates {
    (1, 1.8) (2, 2.0) (3, 2.2) (5, 2.5) (7, 2.7) (10, 2.9)
    (15, 3.1) (20, 3.2) (30, 3.3) (50, 3.35)
};
\addlegendentry{Courbe EIOPA EUR 31/12/2024 (extrapolation)}
\end{axis}
\end{tikzpicture}
\end{center}

\textbf{Publication :} EIOPA publie mensuellement les courbes par devise\\
\textbf{URL :} \url{https://www.eiopa.europa.eu/tools-and-data/risk-free-interest-rate-term-structures}\\
\textbf{Méthode :} Extrapolation Smith-Wilson au-delà de Last Liquid Point (LLP)

\textbf{Référence :} EIOPA Technical Documentation (2020), Règlement Délégué 2015/35 Art. 49
\end{frame}

\begin{frame}{Extrapolation des courbes : méthode Smith-Wilson}
\begin{block}{Problématique industrielle}
\textbf{Question :} Comment actualiser les flux d'un assuré de 30 ans qui vivra potentiellement jusqu'à 100 ans (horizon \highlight{70 ans}) si les marchés ne fournissent des taux liquides que jusqu'à 20-30 ans ?
\end{block}

\begin{columns}[T]
\begin{column}{0.48\textwidth}
\textbf{Last Liquid Point (LLP) :}
\begin{itemize}
    \item Maturité au-delà de laquelle les obligations deviennent illiquides
    \item Pour EUR : LLP ≈ 20 ans
    \item Pour USD : LLP ≈ 30 ans
    \item Données de marché fiables seulement jusqu'au LLP
\end{itemize}

\vspace{0.3cm}
\textbf{Ultimate Forward Rate (UFR) :}
\begin{itemize}
    \item Taux long terme convergence
    \item EUR : UFR = 3.45\% (2025)
    \item Représente anticipations long terme de croissance + inflation
\end{itemize}
\end{column}
\begin{column}{0.48\textwidth}
\textbf{Méthode Smith-Wilson :}
\begin{itemize}
    \item Interpolation/extrapolation lisse
    \item Convergence progressive vers UFR
    \item Convergence complète à 60 ans
    \item Formule mathématique complexe (noyau de Wilson)
\end{itemize}

\vspace{0.3cm}
\begin{center}
\begin{tikzpicture}[scale=0.65]
\begin{axis}[
    xlabel={\scriptsize Maturité},
    ylabel={\scriptsize Taux},
    width=5.5cm,
    height=4.5cm,
    legend style={font=\tiny}
]
\addplot[color=bleuactu, thick] coordinates {
    (1, 1.8) (10, 2.9) (20, 3.2)
};
\addplot[color=rougeactu, thick, dashed] coordinates {
    (20, 3.2) (30, 3.3) (40, 3.4) (60, 3.45) (80, 3.45)
};
\draw[dashed, gray] (axis cs:20,0) -- (axis cs:20,4);
\node[font=\tiny] at (axis cs:20,3.8) {LLP};
\legend{Marché, Extrapolation SW}
\end{axis}
\end{tikzpicture}
\end{center}
\end{column}
\end{columns}

\vspace{0.2cm}
\begin{alertblock}{Impact pratique}
Pour un contrat avec horizon 70 ans, environ \highlight{40-50\% des flux} sont actualisés avec des taux extrapolés (pas observés sur le marché) !
\end{alertblock}

\textbf{Référence :} EIOPA (2020), "Technical documentation of the methodology to derive EIOPA's risk-free interest rate term structures"
\end{frame}

\begin{frame}{Smith-Wilson : formule et paramètres}
\begin{block}{Formule simplifiée}
Le taux forward instantané extrapolé $f(t)$ converge vers l'UFR selon :
\[
f(t) = \text{UFR} + (\text{LLP\_rate} - \text{UFR}) \times e^{-\alpha(t - \text{LLP})} \quad \text{pour } t > \text{LLP}
\]
\end{block}

\textbf{Paramètres clés (EUR, 2025) :}
\begin{table}
\centering
\small
\begin{tabular}{lcc}
\toprule
\textbf{Paramètre} & \textbf{Valeur} & \textbf{Signification} \\
\midrule
\rowcolor{bleuactu!10}
LLP & 20 ans & Dernière maturité liquide \\
UFR & 3.45\% & Taux ultime de convergence \\
\rowcolor{bleuactu!10}
$\alpha$ & 0.10 & Vitesse de convergence \\
Convergence & 60 ans & Maturité où $f(t) \approx$ UFR \\
\bottomrule
\end{tabular}
\end{table}

\begin{exampleblock}{Conséquence pour les assureurs}
\begin{itemize}
    \item Les changements d'UFR par EIOPA impactent massivement le BE Vie (milliards d'euros)
    \item Débats récurrents sur la méthode : taux extrapolés vs taux observés
    \item Certains régulateurs (Suisse) utilisent d'autres méthodes (Nelson-Siegel-Svensson)
\end{itemize}
\end{exampleblock}

\textbf{Point de vigilance :} L'extrapolation introduit un risque modèle significatif en Vie !
\end{frame}

\begin{frame}{Formule complète de projection (récapitulatif)}
\textbf{Algorithme de projection année par année :}

\begin{block}{Étape 1 : Calcul des flux}
\begin{align}
\text{Décès}_t &= N_t \times \overline{\PM}_t \times q_{x+t-1}^d \\
\text{Rachats}_t &= N_t \times \overline{\PM}_t \times q_t^r \\
\text{Frais}_t &= N_t \times \overline{\PM}_t \times f \times (1+\pi)^{t-1} \\
\PB_t &= N_t \times \overline{\PM}_t \times \max(\TMG, 0.90 \times r_{\text{actifs},t}) \\
\text{Primes}_t &= N_t \times P
\end{align}
\end{block}

\begin{block}{Étape 2 : Mise à jour des états}
\begin{align}
N_{t+1} &= N_t \times (1 - q_{x+t-1}^d) \times (1 - q_t^r) \\
\overline{\PM}_{t+1} &= \frac{(\overline{\PM}_t + P - \overline{\PM}_t \times f) \times (1 + i_t)}{(1 - q_{x+t-1}^d)(1 - q_t^r)}
\end{align}
\end{block}

\begin{block}{Étape 3 : Calcul du BE}
\[
\boxed{\BE_0 = \sum_{t=1}^{T} \frac{\text{Décès}_t + \text{Rachats}_t + \text{Frais}_t + \PB_t - \text{Primes}_t}{(1+r_t)^t}}
\]
\end{block}
\end{frame}

% ============================================
% SECTION 4 : INTRODUCTION À PYTHON
% ============================================
\section{Introduction à Python pour l'Actuariat}

\begin{frame}{Limites d'Excel pour l'actuariat industriel}
\begin{alertblock}{Problèmes rencontrés en production}
\begin{enumerate}
    \item \textbf{Volumétrie :} maximum 1 048 576 lignes × 16 384 colonnes
    \begin{itemize}
        \item 50 000 contrats × 40 ans × 100 scénarios = 200 millions de cellules \textcolor{rougeactu}{❌}
    \end{itemize}
    
    \item \textbf{Performance :} calculs matriciels très lents
    \begin{itemize}
        \item Projection 10 000 contrats : 30 min Excel vs \highlight{10 sec} Python
    \end{itemize}
    
    \item \textbf{Maintenance :} formules cachées, erreurs difficiles à détecter
    \begin{itemize}
        \item Risque opérationnel élevé (erreurs de copier-coller, liens cassés)
    \end{itemize}
    
    \item \textbf{Reproductibilité :} difficile à versionner (Git incompatible)
    
    \item \textbf{Collaboration :} conflits de versions, pas de code review possible
    
    \item \textbf{Intégration :} connexion complexe aux bases de données / APIs
\end{enumerate}
\end{alertblock}
\end{frame}

\begin{frame}{Pourquoi Python pour l'actuariat ?}
\begin{columns}[T]
\begin{column}{0.48\textwidth}
\textbf{Avantages de Python :}
\begin{itemize}
    \item \textcolor{vertactu}{Gratuit} et open-source
    \item Écosystème scientifique riche
    \item \textbf{Performance} : vectorisation
    \item Syntaxe claire (pseudo-code)
    \item Intégration facile (SQL, Excel, APIs)
    \item \textbf{Standard de l'industrie}
    \item Traçabilité (Git, tests, CI/CD)
\end{itemize}
\end{column}
\begin{column}{0.48\textwidth}
\textbf{Bibliothèques essentielles :}
\begin{description}
    \item[\texttt{NumPy}] Algèbre linéaire, matrices
    \item[\texttt{Pandas}] DataFrames (tables)
    \item[\texttt{Matplotlib}] Visualisations
    \item[\texttt{SciPy}] Fonctions scientifiques
    \item[\texttt{pyliferisk}] Calculs viagers
    \item[\texttt{SQLAlchemy}] Connexion BDD
\end{description}
\end{column}
\end{columns}

\vspace{0.3cm}
\begin{exampleblock}{Adoption en assurance (2024)}
\begin{itemize}
    \item AXA, Allianz, CNP, Generali : migration Excel → Python en cours
    \item Régulateurs : EIOPA encourage l'utilisation de code reproductible
    \item Cabinets de conseil : Milliman, Deloitte, PWC imposent Python/R
\end{itemize}
\end{exampleblock}
\end{frame}

\begin{frame}[fragile]{Bibliothèque NumPy : calcul vectorisé}
\textbf{Création et manipulation de vecteurs/matrices :}

\begin{lstlisting}
import numpy as np

# Vecteur d'âges
ages = np.arange(25, 66)  # [25, 26, ..., 65]

# Table de mortalité (formule Makeham simplifiée)
qx = 0.0001 * np.exp(0.1 * (ages - 25))

# Probabilités de survie
px = 1 - qx

# Probabilité de survie cumulative jusqu'à chaque âge
survie_cumul = np.cumprod(px)

print(f"Survie de 25 à 65 ans : {survie_cumul[-1]:.2%}")
# Output: Survie de 25 à 65 ans : 98.73%
\end{lstlisting}

\vspace{0.3cm}
\textbf{Avantage majeur :} Opérations vectorisées \highlight{×100 à ×1000 plus rapides} que boucles Python natives !

\textbf{Référence :} \url{https://numpy.org/doc/stable/}
\end{frame}

\begin{frame}[fragile]{Bibliothèque Pandas : manipulation de données tabulaires}
\begin{lstlisting}
import pandas as pd

# Créer un DataFrame (équivalent tableau Excel)
contrats = pd.DataFrame({
    'id': [1, 2, 3, 4, 5],
    'age': [30, 45, 60, 55, 40],
    'PM': [10000, 50000, 80000, 120000, 25000],
    'TMG': [0.01, 0.02, 0.025, 0.015, 0.01]
})

# Filtrer les gros contrats (PM > 50k)
gros_contrats = contrats[contrats['PM'] > 50000]
print(f"Nombre de gros contrats : {len(gros_contrats)}")

# Calculer PM moyenne par tranche d'âge
contrats['tranche'] = pd.cut(contrats['age'], 
                               bins=[0, 40, 60, 100],
                               labels=['<40', '40-60', '>60'])
pm_moyenne = contrats.groupby('tranche')['PM'].mean()
print(pm_moyenne)
# Output:
# <40     17500.0
# 40-60   85000.0
# >60     80000.0
\end{lstlisting}
\end{frame}

\begin{frame}[fragile]{Bibliothèque Matplotlib : visualisations professionnelles}
\begin{lstlisting}
import matplotlib.pyplot as plt
import numpy as np

# Évolution PM avec revalorisation 3%
annees = np.arange(0, 31)
PM = 10000 * (1.03) ** annees

plt.figure(figsize=(10, 6))
plt.plot(annees, PM, marker='o', linewidth=2, 
         color='steelblue', label='PM avec 3% revalorisation')
plt.axhline(y=10000, color='red', linestyle='--', 
            label='PM initiale')
plt.xlabel('Année', fontsize=12)
plt.ylabel('Provision Mathématique (euros)', fontsize=12)
plt.title('Évolution PM sur 30 ans', fontsize=14, fontweight='bold')
plt.legend()
plt.grid(True, alpha=0.3)
plt.tight_layout()
plt.savefig('evolution_PM.png', dpi=300)
plt.show()
\end{lstlisting}

\textbf{Résultat :} Graphique publication-ready exportable en PNG/PDF haute résolution
\end{frame}

\begin{frame}{Exercice Python 1 : Structure de la fonction}
\begin{block}{Objectif}
Créer une fonction \texttt{calculer\_BE\_contrat\_epargne()} qui projette les flux d'un contrat épargne euros et calcule le Best Estimate.
\end{block}

\textbf{Questions à résoudre :}
\begin{enumerate}
    \item Quels paramètres la fonction doit-elle accepter en entrée ?
    \begin{itemize}
        \item Paramètres du contrat (âge, PM, durée, primes)
        \item Hypothèses actuarielles (mortalité, rachat, TMG, frais)
        \item Paramètres économiques (courbe actualisation, inflation)
    \end{itemize}
    
    \item Que doit retourner la fonction ?
    \begin{itemize}
        \item Un DataFrame détaillant les flux année par année
        \item Le Best Estimate total (scalaire)
    \end{itemize}
    
    \item Quelle structure de données utiliser pour stocker les résultats intermédiaires ?
    \begin{itemize}
        \item Liste de dictionnaires ? DataFrame direct ? Tableaux numpy ?
    \end{itemize}
\end{enumerate}

\begin{alertblock}{Astuce}
Pensez à la docstring : documentez clairement les types et significations des paramètres !
\end{alertblock}
\end{frame}

\begin{frame}{Exercice Python 2 : Boucle de projection}
\begin{block}{Contexte}
Vous devez maintenant implémenter la boucle qui projette année par année les flux futurs.
\end{block}

\textbf{Questions guidées :}
\begin{enumerate}
    \item \textbf{Initialisation :} Quelles variables initialisez-vous avant la boucle ?
    \begin{itemize}
        \item Liste vide pour stocker résultats
        \item PM initiale, nombre de têtes initial (= 1.0)
    \end{itemize}
    
    \item \textbf{À chaque itération $t$ :}
    \begin{itemize}
        \item Comment calculer l'âge de l'assuré en année $t$ ?
        \item Comment récupérer $q_x$ dans la table de mortalité (dictionnaire) ?
        \item Comment récupérer $q^r_t$ dans le tableau numpy des rachats ?
    \end{itemize}
    
    \item \textbf{Calcul des flux :}
    \begin{itemize}
        \item Décès = $PM \times q_x^d \times N_t$ (formule vue en cours)
        \item Rachats = $PM \times q_t^r \times N_t$
        \item Frais proportionnels = $PM \times f \times N_t \times (1 + \pi)^{t-1}$
        \item \highlight{Frais unitaires} = $F_0 \times N_t \times (1 + \pi)^{t-1}$ (ex: 30€/contrat)
        \item Frais totaux = frais proportionnels + frais unitaires
        \item Primes = $P \times N_t$
    \end{itemize}
\end{enumerate}
\end{frame}

\begin{frame}{Exercice Python 3 : Actualisation et mise à jour}
\textbf{Questions (suite) :}

\begin{enumerate}
    \setcounter{enumi}{3}
    \item \textbf{Actualisation :}
    \begin{itemize}
        \item Flux net = flux sortants - flux entrants
        \item Facteur actualisation = $\frac{1}{(1+r_t)^t}$ où $r_t$ est le taux à maturité $t$
        \item Flux actualisé = flux net × facteur actualisation
    \end{itemize}
    
    \item \textbf{Mise à jour des états :}
    \begin{itemize}
        \item Nouvelle PM : formule vue en cours avec revalorisation TMG
        \item Attention : calculer PM par tête survivante !
        \item Nouveau nombre de têtes : $N_{t+1} = N_t \times (1-q_x^d) \times (1-q_t^r)$
    \end{itemize}
    
    \item \textbf{Stockage :}
    \begin{itemize}
        \item Créer un dictionnaire avec tous les résultats de l'année
        \item L'ajouter à la liste résultats avec \texttt{.append()}
    \end{itemize}
\end{enumerate}

\begin{exampleblock}{Piège classique}
Ne pas oublier de tester \texttt{if nb\_tetes > 0} avant de diviser par nb\_tetes (éviter division par zéro) !
\end{exampleblock}
\end{frame}

\begin{frame}{Exercice Python 4 : Finalisation et utilisation}
\textbf{Questions finales :}

\begin{enumerate}
    \item \textbf{Après la boucle :}
    \begin{itemize}
        \item Comment convertir la liste de dictionnaires en DataFrame pandas ?
        \item Comment calculer le BE total à partir de la colonne des flux actualisés ?
    \end{itemize}
    
    \item \textbf{Préparation des données d'entrée :}
    \begin{itemize}
        \item Comment créer un dictionnaire \texttt{table\_mortalite} avec formule $q_x = 0.0001 \times e^{0.08(x-18)}$ ?
        \item Comment créer un tableau numpy pour la loi de rachat : [15\%, 15\%, 10\%, 10\%, 6\%, 6\%, 6\%, 6\%, 3\%, ...] ?
        \item Comment créer une courbe flat à 2.5\% sur 40 ans avec numpy ?
    \end{itemize}
    
    \item \textbf{Appel de la fonction :}
    \begin{itemize}
        \item Avec quels paramètres tester ? (exemple : homme 30 ans, PM = 100 000 €)
        \item Comment afficher le résultat proprement ?
    \end{itemize}
\end{enumerate}

\begin{alertblock}{Challenge supplémentaire 1 : Frais unitaires}
Modifier votre fonction pour ajouter un paramètre \texttt{frais\_unitaires=30} et intégrer les frais fixes de 30€/contrat/an indexés sur l'inflation. Comparer le BE avec et sans frais unitaires !
\end{alertblock}

\begin{alertblock}{Challenge supplémentaire 2 : Visualisation}
Créer un graphique matplotlib montrant l'évolution de la PM et du nombre de têtes sur deux axes !
\end{alertblock}
\end{frame}

\begin{frame}{Comparaison Excel vs Python : tableau de synthèse}
\begin{table}
\centering
\small
\begin{tabular}{l|c|c}
\toprule
\textbf{Critère} & \textbf{Excel} & \textbf{Python} \\
\midrule
\rowcolor{bleuactu!10}
Lisibilité du code & ⭐⭐ & ⭐⭐⭐⭐⭐ \\
Maintenance & ⭐⭐ & ⭐⭐⭐⭐⭐ \\
\rowcolor{bleuactu!10}
Performance (vitesse) & ⭐⭐ & ⭐⭐⭐⭐⭐ \\
Scalabilité (volumétrie) & ⭐ & ⭐⭐⭐⭐⭐ \\
\rowcolor{bleuactu!10}
Courbe d'apprentissage & ⭐⭐⭐⭐⭐ & ⭐⭐⭐ \\
Visualisation rapide & ⭐⭐⭐⭐⭐ & ⭐⭐⭐⭐ \\
\rowcolor{bleuactu!10}
Traçabilité (Git) & ⭐ & ⭐⭐⭐⭐⭐ \\
Intégration (BDD, APIs) & ⭐⭐ & ⭐⭐⭐⭐⭐ \\
\rowcolor{bleuactu!10}
Coût & ⭐⭐⭐ (licence Office) & ⭐⭐⭐⭐⭐ (gratuit) \\
\bottomrule
\end{tabular}
\end{table}

\vspace{0.3cm}
\begin{block}{Recommandation pratique}
\begin{itemize}
    \item \textbf{Excel :} prototypes, analyses exploratoires, 1-100 contrats, présentations rapides
    \item \textbf{Python :} production industrielle, >100 contrats, scénarios multiples, reporting ACPR
\end{itemize}
\end{block}
\end{frame}

% ============================================
% SECTION 5 : SYNTHÈSE ET PERSPECTIVES
% ============================================
\section{Synthèse et Perspectives}

\begin{frame}{Récapitulatif des concepts clés}
\begin{block}{Compétences acquises dans cette séance}
\begin{enumerate}
    \item \textbf{Typologie Vie :} Distinguer euros, UC, prévoyance, rentes viagères
    \item \textbf{Modélisation :} Projeter flux futurs (prestations, primes, frais, PB)
    \item \textbf{Hypothèses :} Appliquer tables mortalité, lois de rachat, inflation frais
    \item \textbf{Méthodes :} Choisir entre police-par-police et model points
    \item \textbf{Actualisation :} Utiliser courbe EIOPA sans risque
    \item \textbf{Best Estimate :} Calculer le BE conforme Solvency II
    \item \textbf{Programmation :} Coder en Python pour industrialiser
\end{enumerate}
\end{block}

\vspace{0.3cm}
\begin{alertblock}{Compétences opérationnelles}
\begin{itemize}
    \item[$\checkmark$] Projeter un portefeuille vie sur Excel (petite échelle)
    \item[$\checkmark$] Écrire un script Python de projection (industrialisable)
    \item[$\checkmark$] Réaliser des analyses de sensibilité (TMG, rachats, mortalité)
    \item[$\checkmark$] Interpréter les résultats et identifier les drivers du BE
\end{itemize}
\end{alertblock}
\end{frame}

\begin{frame}{Limites de l'approche déterministe}
\begin{alertblock}{Ce cours = simplification pédagogique}
En production réelle pour Solvency II, les compagnies utilisent des méthodes plus avancées :
\end{alertblock}

\begin{block}{1. Projections stochastiques (Monte-Carlo)}
\begin{itemize}
    \item Taux d'intérêt modélisés (CIR, Hull-White, Black-Karasinski)
    \item 1 000 à 10 000 scénarios économiques (ESG)
    \item Valorisation options/garanties par Monte-Carlo
\end{itemize}
\end{block}

\begin{block}{2. Rachats dynamiques comportementaux}
\begin{itemize}
    \item Modèles économétriques calibrés sur historique
    \item Dépendance au contexte économique (spread taux)
    \item Stress massif de rachats (Solvency II SCR)
\end{itemize}
\end{block}

\begin{block}{3. Modélisation actif-passif jointe (ALM)}
\begin{itemize}
    \item Projection conjointe bilan actif et passif
    \item Gestion dynamique allocation d'actifs
    \item Calcul du SCR (capital de solvabilité requis)
\end{itemize}
\end{block}

\textbf{Ces techniques seront vues en Master 2 Actuariat}
\end{frame}

\begin{frame}{Niveaux de complexité : du cours à la production}
\begin{center}
\begin{tikzpicture}[scale=1.0]
% Axes
\draw[->] (0,0) -- (11,0) node[right] {\scriptsize Complexité};
\draw[->] (0,0) -- (0,6) node[above] {\scriptsize Temps calcul};

% Points
\node[circle, fill=vertactu, minimum size=0.4cm, label=below:{\scriptsize Excel simple}] (excel) at (2, 1) {};
\node[circle, fill=bleuactu, minimum size=0.4cm, label=below:{\scriptsize Python déterministe}] (python) at (5, 2) {};
\node[circle, fill=orangeactu, minimum size=0.4cm, label=below:{\scriptsize Python stochastique}] (stoch) at (8, 4) {};
\node[circle, fill=rougeactu, minimum size=0.4cm, label=below:{\scriptsize Production ALM}] (prod) at (10, 5.5) {};

% Flèches progression
\draw[->, thick, dashed] (excel) -- (python);
\draw[->, thick, dashed] (python) -- (stoch);
\draw[->, thick, dashed] (stoch) -- (prod);

% Annotations
\node[align=center, font=\tiny] at (2, 2.5) {1 contrat\\10 min\\Séance 4};
\node[align=center, font=\tiny] at (5, 3.5) {100 contrats\\5 min\\Séance 4};
\node[align=center, font=\tiny] at (8, 5.5) {1000 scénarios\\1 h\\Master 2};
\node[align=center, font=\tiny] at (10, 6.8) {Full ALM\\12 h\\Production};
\end{tikzpicture}
\end{center}

\vspace{0.3cm}
\textbf{Votre progression :}
\begin{itemize}
    \item \textcolor{vertactu}{\textbf{Étape 1-2 (ce cours) :}} maîtrise Excel et Python déterministe
    \item \textcolor{orangeactu}{\textbf{Étape 3 (M2, stage) :}} méthodes stochastiques avancées
    \item \textcolor{rougeactu}{\textbf{Étape 4 (carrière) :}} production industrielle full ALM
\end{itemize}
\end{frame}

\begin{frame}{Pour aller plus loin : ressources}
\begin{columns}[T]
\begin{column}{0.48\textwidth}
\textbf{Réglementation :}
\begin{itemize}
    \item Directive 2009/138/CE (Solvency II)
    \item Règlement Délégué 2015/35
    \item EIOPA Guidelines on Technical Provisions (2015)
    \item ACPR Position 2013-P-01
    \item QIS5 Technical Specifications
\end{itemize}

\vspace{0.3cm}
\textbf{Ouvrages techniques :}
\begin{itemize}
    \item Planchet \& Thérond, \textit{Modélisation statistique des phénomènes de durée}
    \item Devineau \& Loisel, \textit{Risque de taux en assurance-vie}
    \item Wüthrich \& Merz, \textit{Stochastic Claims Reserving}
\end{itemize}
\end{column}
\begin{column}{0.48\textwidth}
\textbf{Python pour actuaires :}
\begin{itemize}
    \item Doc NumPy/Pandas/Matplotlib
    \item Bibliothèque \texttt{pyliferisk}
    \item Bibliothèque \texttt{actuarialmath}
    \item Cours Coursera/DataCamp
\end{itemize}

\vspace{0.3cm}
\textbf{Sites de référence :}
\begin{itemize}
    \item Institut des Actuaires (France)
    \item Society of Actuaries (USA)
    \item Actuarial Outpost (forum)
    \item EIOPA website (données/guidelines)
\end{itemize}
\end{column}
\end{columns}

\vspace{0.3cm}
\begin{block}{Formation continue}
\textbf{Veille réglementaire :} Bulletin ACPR, EIOPA Opinions, Guidelines mises à jour\\
\textbf{Veille technique :} Conférences IA, nouveaux packages Python, GPU computing
\end{block}
\end{frame}

\begin{frame}{De l'étudiant à l'actuaire opérationnel}
\begin{block}{Dans un service actuariat production d'une compagnie :}
\begin{multicols}{2}
\begin{itemize}
    \item Portefeuilles 100 000 - 500 000 contrats
    \item Projections sur clusters HPC (10-100 cœurs)
    \item Temps de calcul : plusieurs heures
    \item Code Python versionné (Git)
    \item Tests unitaires, CI/CD
    \item Peer review obligatoire
    \item BDD SQL (PostgreSQL, Oracle)
    \item Dashboards (Tableau, Power BI)
    \item Reporting ACPR/EIOPA (QRT)
\end{itemize}
\end{multicols}
\end{block}

\vspace{0.3cm}
\begin{exampleblock}{Le chemin pour y arriver}
\begin{enumerate}
    \item[$\checkmark$] \textbf{Concepts théoriques :} maîtrisés (ce cours)
    \item[$\checkmark$] \textbf{Implémentation simple :} Excel et Python basique (TD)
    \item[$\rightarrow$] \textbf{Projets complexes :} stages, alternance, mémoire actuariat
    \item[$\rightarrow$] \textbf{Bonnes pratiques dev :} Git, tests, documentation, review
    \item[$\rightarrow$] \textbf{Formation continue :} veille réglementaire, nouveaux outils
\end{enumerate}
\end{exampleblock}

\textbf{Vous êtes maintenant équipés pour devenir des actuaires opérationnels !}
\end{frame}

\begin{frame}[c]  % [standout] n'existe que dans le thème metropolis
\Huge Questions ?

\vspace{1.5cm}

\Large
\textbf{Prochaine séance :}\\
Risk Margin et SCR Vie

\vspace{1cm}

\normalsize
\textit{« Un actuaire n'est pas quelqu'un qui aime les mathématiques, \\
mais quelqu'un qui aime résoudre des problèmes avec des mathématiques »}

\vspace{1.5cm}

Contact : \texttt{fabrice.borel-mathurin@acpr.banque-france.fr}
\end{frame}

\end{document}
